\title{Byte Latent Transformer: Patches Scale Better Than Tokens}

\begin{abstract}
We present the Byte Latent Transformer ({\textsc BLT}), a novel architecture for large language models operating at the byte level, which for the first time achieves parity with traditional tokenization-based LLMs in terms of performance at scale while substantially advancing both inference speed and robustness. 
{\textsc BLT} transforms byte sequences into adaptive-length patches, designating these as the fundamental computational unit. The process segments patches guided by the entropy of the following byte, focusing computational power and model attention on data regions exhibiting higher complexity. We conduct the first scaling study with flop-constrained training for byte-level models, reaching up to 8B parameters and training with 4T bytes of data. Our findings confirm that effective scaling is possible when training on raw bytes without reliance on a static vocabulary. Efficiency in training and inference improves as the system dynamically opts for extended patches where data exhibits predictability, producing further gains in tasks requiring reasoning and handling of rare phenomena. Summarizing our findings, for any given inference budget, {\textsc BLT} attains superior scaling compared to tokenization-based counterparts by enabling simultaneous increases in both the model and patch size.
\end{abstract}

\section{Introduction}
\label{section:intro}

We present the Byte Latent Transformer~(\textbf{BLT}), an architecture that dispenses with tokenization and instead learns directly from raw byte sequences. BLT achieves, for the first time, parity with models relying on tokenization-based pipelines at scale, while delivering marked gains in both computational efficiency and robustness (\S\ref{sec:robustness}).

Contemporary large language models (llms) generally operate in an almost fully end-to-end manner, with the notable exception of a tokenization stage—this heuristic preprocessing aggregates bytes into predetermined token sets. This process introduces compression biases within strings, resulting in problems such as heightened susceptibility to specific domains/modalities~\citep{dagangetting}, increased vulnerability to input noise~(\S\ref{sec:robustness}), a deficiency in orthographic generalization~\citep{edman2024cute}, and inequitable treatment across languages~\citep{liang2023xlm,petrov2024language,limisiewicz2024myte}.

Tokenization has historically played a crucial role since training llms directly on bytes incurs excessive computational expense, largely due to the extended sequence lengths required~\citep{xue2022byt5}. Previous research addresses this by introducing more computationally efficient self-attention mechanisms~\citep{el2020characterbert, clark2022canine}, or by designing architectures that eliminate attention altogether~\citep{wang2024mambabyte}~(\S\ref{section:related_work}). Nonetheless, such advancements predominantly benefit \textit{small models}. When considering large-scale models, most of the computational burden within a Transformer arises from the sizable feed-forward network layers that process each byte, rather than from the attention component.

For optimal compute distribution, we introduce an adaptable and trainable technique for assembling bytes into \textit{patches}~(\S\ref{section:patching}), alongside a novel architecture that integrates both byte-level and patch-level data. In contrast to tokenization, {\textsc BLT} does not employ a predetermined patch vocabulary. Instead, flexible sequences of bytes are transformed into latent patch embeddings through compact, parameterized encoder and decoder components. Our empirical findings demonstrate that this strategy allocates computational resources \textit{more} efficiently than approaches relying on tokenization.

Conventional tokenization-based LLMs distribute computational resources equally across all tokens, which prioritizes uniform efficiency but does not account for varying prediction complexities. The compression strategies employed in token induction may not accurately reflect the actual cognitive load required for each prediction. Our approach departs from this uniform allocation by advocating for adaptive compute distribution, focusing resources where prediction difficulty is greater. For instance, generating the suffix of a word often represents a straightforward, low-entropy choice, and thus does not necessitate the extensive capacity of a full transformer model; in contrast, predicting a sentence’s initial word involves a more challenging, high-entropy process. This principle guides the design of {\textsc BLT}'s architecture~(\S\ref{section:architecture}), which features a trio of transformer modules: two compact, byte-level \textit{local models} and a more substantial \textit{latent transformer} operating at a global scale~(\autoref{fig:arch_high_level}). 

To efficiently assign computational effort, {\textsc BLT} employs a strategy in which bytes are organized into patches according to the entropy associated with predicting the subsequent byte. This results in context-dependent groupings of bytes, each exhibiting a similar density of informational content.

We introduce the first scaling analysis of byte-level models governed by fixed computational budgets, exploring models with up to 8B parameters and 4T bytes of training data. Our findings demonstrate the feasibility of training models end-to-end at scale using raw bytes, thereby eliminating the need for predetermined vocabulary tokenization. In terms of training efficiency, {\textsc BLT} achieves parity with Llama 3 under equivalent flop-controlled conditions\footnote{Our measure of computational expense is based on the total Floating Point OPerations (flops) required.}, while outperforming it by consuming up to 50\% fewer flops at inference~(\S\ref{section:scaling}). 

Furthermore, we find that modeling raw bytes directly leads to notable gains when handling rare and diverse patterns within the data. The {\textsc BLT} models exhibit greater robustness compared to their tokenizer-reliant counterparts, particularly when presented with noisy inputs, and show superior ability at the character level across orthographic recognition, phonological reasoning, and tasks involving low-resource machine translation~(\S\ref{sec:robustness}). 

Additionally, with {\textsc BLT}, it is possible to concurrently expand both model size and input patch size without increasing the inference flop budget. Utilizing longer patch sizes generally decreases the frequency of global latent transformer invocations, conserving computational resources that can instead be directed to enlarging the model’s latent capacity. Our inference flop-constrained scaling studies~(\autoref{fig:fixed_inference_scaling}) reveal that this approach produces markedly improved scaling behaviors relative to models built on tokenization-based strategies.

To conclude, this work provides the following key advances: 1) We present {\textsc BLT}, a byte latent llm framework that adaptively distributes computational resources for increased flop efficiency, 2) We demonstrate parity with Llama 3 in terms of training flops up to 8B parameters, with the added flexibility to sacrifice minor performance on evaluation metrics in exchange for up to 50\% greater flop efficiency, 3) {\textsc BLT} enables a novel scaling paradigm for llms, allowing expansion of model size while adhering to a fixed inference compute budget, 4) Our results highlight the superior robustness of {\textsc BLT} models against noisy input and their capacity to capture sub-word information often overlooked by token-based llms. The codebase for both training and inference with {\textsc BLT} is made publicly available at~\url{https://github.com/facebookresearch/blt}.

\section{Patching: From Individual Bytes to Groups of Bytes}
\label{section:patching}

Dividing bytes into discrete \textit{patches} enables {\textsc BLT} to assign computational resources adaptively, depending on contextual demands. Figure~\ref{fig:patching_types} illustrates various approaches for segmenting byte sequences into patches. More formally, a patching function $f_p$ takes a sequence of bytes $\pmb{x} = \{x_i | i = 1, \ldots, n\}$ of length $n$, and partitions it into a sequence of $m < n$ patches $\pmb{p} = \{p_j | j = 1, \ldots, m\}$ by assigning each $x_i$ to the set \{0, 1\}, where a value of 1 signifies the beginning of a new patch. For models that operate on tokens or patches, the primary factor dictating compute usage is the number of processing steps performed by the main Transformer; in {\textsc BLT}, this is equivalent to the number of patches required by a specific patching function to represent the input data. As a result, the expected length of a patch, referred to as \textit{patch size}, becomes the central determinant of compute cost during both the training and inference phases when utilizing a particular patching method~(\S\ref{section:flops}). 

In the following sections, we define three distinct patching approaches: fixed-length patching~(\S\ref{section:static-patch}), patching at whitespace boundaries~(\S\ref{section:white-patch}), and a dynamic approach leveraging entropy estimates from a lightweight byte language model~(\S\ref{section:dyn-patch}). We conclude with an overview of incremental patching strategies and an exploration of how patching contrasts with tokenization~(\S\ref{section:bpe-patch}).

\subsection{Strided Patching Every K Bytes}
\label{section:static-patch}

A commonly used approach for grouping bytes involves dividing them into fixed-size patches of length $k$, as employed in MegaByte~\citep{yu2023megabyte}. Utilizing a constant stride simplifies both implementation during training and inference, while also providing an easy way to adjust the typical patch size and, consequently, the computational load. Despite these advantages, fixed patching presents notable drawbacks. Primarily, computational resources are not adaptively directed to regions where they are most essential: a transformer step $j$ might be underutilized when predicting less informative tokens like whitespace, or, conversely, might not provide adequate computation for highly informative regions such as complex mathematical notation. Additionally, this method results in arbitrary and non-context-aware segmentation of byte sequences—for instance, leading to inconsistent splitting of identical words.

\subsection{Space Patching}
\label{section:white-patch}

\citet{slagle2024spacebyte} introduces a straightforward and effective modification to strided patching by generating new patches immediately after any space-like byte\footnote{
    Space-like bytes are specified as bytes other than Latin letters, digits, or UTF-8 continuation bytes; furthermore, every patch must include at least one byte that is not space-like.}, leveraging these boundaries as natural demarcations for linguistic segments in various languages. Under Space patching, each word is allotted a step of the latent transformer (which requires greater computational effort), thereby standardizing the patching of words throughout sequences and concentrating computational resources on more challenging prediction tasks that usually succeed spaces. To illustrate, the challenge of predicting the initial byte in response to ``Who composed the Magic Flute?\uline{\hspace{.6em}}'' is significantly greater than that of predicting the subsequent bytes after the ``M,'' as the opening character narrows the spectrum of possible continuations, leaving the completion of ``Mozart'' more predictable. Nevertheless, space patching is not universally applicable across all languages and domains, and its primary limitation is its fixed patch size. In the following section, we describe a new patching strategy inspired by the observation that the initial bytes of words generally pose the greatest predictive difficulty, and that also enables flexible adjustment of patch size.

\subsection{Entropy Patching: Using Next-Byte Entropies from a Small Byte LM}
\label{section:dyn-patch}

Instead of utilizing rule-based heuristics like whitespace to determine boundaries, we adopt a data-centric method focused on detecting regions with elevated prediction uncertainty for the next byte. We propose \textit{entropy patching}, a technique that leverages entropy metrics to establish the patch boundaries.

To this end, we train a compact auto-regressive language model operating at the byte level on the {\textsc BLT} training set and evaluate next-byte entropy according to the language model distribution $p_{e}$ over the byte vocabulary $\mathcal{V}$:

\begin{equation}
H(x_i) = \sum_{v \in \mathcal{V}} p_{e}(x_i=v|\pmb{x}_{<i}) \log p_{e}(x_i=v|\pmb{x}_{<i})
\end{equation}

We explore two strategies for detecting patch boundaries using entropies $H(x_i)$. The first strategy selects points exceeding a universal entropy threshold, as depicted in~\autoref{fig:patching}. The second strategy focuses on points where the entropy is significantly greater than that of the preceding point. This latter approach effectively captures locations where the entropy pattern departs from an approximately monotonically decreasing trend within the patch.

\begin{align*}
    \text{Global Constraint}\hspace{1cm}H(x_t)& > \theta_g\\
    \text{Approx. Monotonic Constraint}\hspace{1cm}H(x_t)& - H(x_{t-1}) > \theta_r
\end{align*}

During the data loading phase, a lightweight preprocessing procedure is utilized to detect patch boundaries. This contrasts with the approach of~\citet{nawrot-etal-2023-efficient}, where a classifier is employed to learn and predict patch boundaries based on entropy measures. In our experiments~(\S\ref{section:experiments}), we evaluate and contrast these two techniques for differentiating between bytes of low and high entropy.

\subsection{The Byte-Pair Encoding (BPE) Tokenizer and Incremental Patching}
\label{section:incremental-patching}
\label{section:bpe-patch}

Numerous contemporary llms, such as our baseline model Llama 3, implement a subword tokenization strategy, notably bpe~\citep{gage1994new, sennrich-etal-2016-neural}. In this work, the term ``tokens'' denotes clusters of bytes selected from a \textit{finite} set established before the training process. In contrast, ``patches'' are dynamically created groupings of sequences that lack a predetermined vocabulary. One key distinction is that, when utilizing tokens, the model is unable to directly observe or manipulate the underlying byte-level features.

An important advancement introduced by {\textsc BLT} in comparison to tokenization-based models lies in its reconfiguration of the balance between vocabulary size and computational demands. In conventional llms, enlarging the vocabulary typically results in tokens that, on average, are longer. This reduces the total number of decoding steps required, yet simultaneously expands the output dimension of the model’s final projection layer, increasing computational load. This inherent compromise restricts tokenization approaches from meaningfully diversifying token length or reducing inference time. As an illustrative case, Llama 3 achieves a rise in mean token length from 3.7 to 4.4 bytes, but this comes at the expense of quadrupling its embedding table size in comparison to Llama 2.

During generation, {\textsc BLT} must determine at each step in the byte sequence whether it is positioned at a patch boundary, since this dictates when the Latent Transformer is called for further computation. This decision-making process must proceed without any reliance on future, as yet ungenerated, parts of the sequence. Consequently, the patching method must operate solely on already generated bytes, with no foresight into upcoming bytes, to determine how the sequence should be divided. More formally, an incremental patching scheme $f_p$ is defined by the following property:
$$f_p(\pmb{x}_{<i}) = f_p(\pmb{x})_{<i}$$
The commonly used bpe tokenization does not qualify as an incremental patching scheme because the same byte prefix may be segmented in varying ways depending on the subsequent content, and thus fails to meet this requirement\footnote{Inserting a unique delimiter token at patch boundaries can transform bpe into an incremental scheme, albeit at the expense of elongating the byte sequence.}.

\section{{\textsc BLT} Architecture}
\label{section:architecture}
{\textsc BLT} comprises a large, globally acting autoregressive language model that processes patch-level representations. This component is supported by two compact local models: one to convert byte sequences into patches and another to transform patch representations back into byte sequences (\autoref{fig:arch_high_level}).

\subsection{Latent Global Transformer Model}

\textit{The Latent Global Transformer} refers to an autoregressive transformer model, denoted as $\mathcal{G}$ with $l_{\mathcal{G}}$ layers, that transforms a sequence of latent input patch representations $p_j$ into corresponding output patch representations $o_j$. In this work, the index $j$ is consistently assigned to patches and $i$ to individual bytes. The architecture implements a block-causal attention mask~\citep{dubey2024llama}, thereby confining each position’s attention span to patches up to and including the current patch within a given document. Since the global model accounts for the majority of floating point operations during both pre-training and inference phases, judicious selection of when to deploy this component enables adaptive allocation of computational resources in proportion to the complexity of specific input and output regions.

\subsection{Local Encoder}
\label{section:local}

\textit{The Local Encoder Model}, represented by $\mathcal{E}$, utilizes a compact transformer-based structure characterized by $l_{\mathcal{E}} << l_{\mathcal{G}}$ layers. Its core function is to transform input byte sequences $b_i$ into rich patch representations $p_j$ in an efficient manner. Distinguishing itself from the conventional transformer setup, this model incorporates a cross-attention layer immediately following each transformer layer; this mechanism aggregates the byte-level representations into higher-level patch representations~(\autoref{fig:crossattn}).

To begin, each input byte $b_i$ is projected into an embedding space via a matrix of size $\mathbb{R}^{256 \times h_{\mathcal{E}}}$, yielding $x_i$ embeddings. These embeddings may be further supplemented with additional signals using hash-embeddings~(\S\ref{sec:hash-emb}) if desired. Subsequently, the sequence is processed through alternating stacks of transformer and cross-attention modules~(\S\ref{sec:enc-cross-attn}) which convert the byte-level embeddings into patch representations $p_i$ suited for consumption by the global transformer $\mathcal{G}$.

The attention mechanism within each transformer layer leverages a \textit{local block causal} masking strategy, confining each byte's attention to a sliding window of $w_{\mathcal{E}}$ prior bytes. This window can extend across adaptive patch boundaries, though it does not span across document demarcations. The specifics concerning the embedding scheme and the structure of the cross-attention component are elaborated in the subsections that follow.

\subsubsection{Encoder Hash n-gram Embeddings}
\label{sec:ngram-emb}
\label{sec:hash-emb}
An essential aspect of building resilient and informative representations at each position $i$ involves encoding the context provided by previous bytes. Within {\textsc BLT}, this is accomplished by representing not only the individual byte $b_i$ but also its participation in a sequence of bytes, forming a byte n-gram. At every step $i$, we begin by assembling byte-grams

\begin{align}
    g_{i,n}=\{b_{i-n + 1},\ldots, b_i\}
\end{align}

for each byte position $i$ and $n$ from three to eight.\footnote{
We omit byte-grams of size $n$ or more when $i<n$. }

Next, we present hash $n$-gram embeddings, which utilize a hash function to assign each byte $n$-gram to a position within a fixed-size embedding table $E_{n}^{hash}$ for every $n \in\{3,4,5,6,7,8\}$~\citep{bai2010learning}. This hash-based $n$-gram embedding is combined with the corresponding byte embedding, after which the sum is normalized and forwarded as input to the local encoder module. The computation for the augmented embedding is

\begin{align}
e_i &= x_i + \sum_{n = 3,...,8}E_{n}^{hash}(\text{Hash}(g_{i,n})) \\
\text{where, Hash}(g_{i,n}) &= \text{RollPolyHash}(g_{i,n}) \% |E_{n}^{hash}|
\end{align}

We divide $e_i$ by the total number of $n$-gram sizes incremented by one, and apply RollPolyHash as outlined in Appendix~\ref{appendix:rollpolyhash}. Section~\ref{section:ablations} presents an analysis of how varying $n$ in $n$-gram hash embeddings, as well as changing the embedding table size, influences flop-controlled scaling law behavior. Beyond using hash-based $n$-gram embeddings, we also investigated embeddings determined by $n$-gram frequency, with further details on these experiments discussed in Appendix~\ref{appendix:ngrams}.

\subsubsection{Encoder Multi-Headed Cross-Attention}
\label{sec:enc-cross-attn}
Our design largely mirrors the input cross-attention mechanism outlined in the Perceiver framework~\citep{jaegle2021perceiver}. The key distinction in our approach is that the latent states here represent variable-length patch encodings instead of being drawn from a static collection of latent vectors~(\autoref{fig:crossattn}). Moreover, each latent state exclusively attends to the bytes constituting its own patch. This module is built around a query vector for each patch $p_j$, initialized via an aggregation (such as pooling) of the byte-level representations pertaining to $p_j$. This is subsequently passed through a linear transformation, $\mathcal{E}_{C} \in \mathbb{R}^{h_{\mathcal{E}} \times (h_{\mathcal{E}} \times U_{\mathcal{E}})}$, with $U_{\mathcal{E}}$ denoting the count of cross-attention heads in the encoder. More precisely, defining $f_{\text{bytes}}(p_j)$ as the ordered set of bytes from patch $p_j$, the cross-attention computation can be expressed as

\begin{align}
P_{0,j} &= \mathcal{E}_{C}(f_\text{bytes}(p_j)), \quad f~ \text{denotes a pooling operation} \\
P_l &= P_{l-1} + W_o\left(\text{softmax}\left(\frac{QK^T}{\sqrt{d_k}}\right)V\right) \\
\text{where } Q_j &= W_q(P_{l-1,j}), \quad K_i = W_k(h_{l-1,i}), \quad V_i = W_v(h_{l-1,i}) \\
h_l &= \text{Encoder-Transformer-Layer}_l(h_{l-1})
\end{align}

Here, $P \in \mathbb{R}^{n_p \times h_{\mathcal{G}}}$ denotes the set of $n_p$ patch representations that are to be handled by the global model; this matrix is formed by aggregating the byte embeddings $e_i$ associated with each patch $p_j$. The projection matrices $W_q$, $W_k$, $W_v$, and $W_o$ correspond, respectively, to queries, keys, values, and outputs, with the keys and values derived from linearly projecting the byte representations $h_i$ generated by the previous layer ($e_i$ in the initial layer). Our approach incorporates a tailored masking mechanism such that every query $Q_j$ can attend exclusively to the keys and values that are linked to bytes within patch $j$. Because multi-head attention is applied to $Q, K$, and $V$, and because patch representations usually possess a higher dimension ($h_{\mathcal{G}}$) relative to $h_{\mathcal{E}}$, we represent $P_l$ as multiple attention heads each of size $h_{\mathcal{E}}$ during cross-attention, concatenating these heads afterward to recover the full $h_{\mathcal{G}}$ dimensionality. Further, we apply a pre-LayerNorm step to the queries, keys, and values, and omit positional encodings within this cross-attention mechanism. Lastly, a residual connection is included to wrap around the cross-attention layer.

\subsection{Local Decoder} 

The local decoder $\mathcal{D}$, much like its encoder counterpart, employs a compact transformer-based architecture comprising $l_{\mathcal{D}}$ layers, where $l_{\mathcal{D}}$ is significantly smaller than $l_{\mathcal{G}}$. This module converts a series of global patch embeddings $o_j$ into the output raw bytes $y_i$. In predicting each byte in the sequence, the local decoder conditions on the sequence of bytes already produced, utilizing the hidden states generated by the local encoder for the input byte sequence. The architecture alternates between $l_{\mathcal{D}}$ cross-attention layers and transformer blocks. Importantly, in each layer, the cross-attention mechanism is positioned before the transformer block, mapping patch-level representations to byte-level representations. The subsequent transformer layer then processes these byte representations.

\subsubsection{Decoder Multi-headed Cross-Attention} 

Within the decoder's cross-attention mechanism, the function of queries and key/values is reversed: byte-representations act as queries, while patch representations serve as keys and values. The initial set of byte-representations employed for cross-attention is seeded using the byte embeddings derived from the output of the encoder's final layer, i.e., $h_{l_{\mathcal{E}}}$. The byte-representations for the $l$-th layer, denoted $d_{l,i}$, are obtained by:

\begin{align}
D_0 &= h_{l_{\mathcal{E}}} \\
B_l &= D_{l-1} + W_o\left(\text{softmax}\left(\frac{QK^T}{\sqrt{d_k}}\right)V\right), \\
  &\text{where } Q_i = W_q(d_{l-1,i}), K_i = W_k(\mathcal{D}_C(o_j)), V_i = W_v(\mathcal{D}_C(o_j)) \\
  D_l &= \text{Decoder-Transformer-layer}_l(B_l)
\end{align}

Here, $W_k$ and $W_v$ denote the projection matrices for keys and values, respectively, each performing a linear mapping along with a splitting operation $\mathcal{D}_C$, which are applied to the final patch representations $o_j$ yielded by the global model. $W_q$ serves as the query projection matrix acting on the byte representations $d_{l-1}$ from the previous layer of the decoder transformer, or on $h_{l_{\mathcal{E}}}$ in the case of the initial layer. $W_o$ represents the output projection matrix, making $B$ a matrix in $\mathbb{R}^{h_{\mathcal{D}} \times n_b}$, where $n_b$ denotes the count of output bytes. The subsequent decoder representations $D_l$ result from processing the output $B$ of the cross-attention mechanism through a transformer layer in the decoder. As with the cross-attention in the local encoder, the architecture employs multiple attention heads, applies pre LayerNorm normalization, excludes positional embeddings, and includes a residual connection surrounding the cross-attention component.

\section{Experimental Setup}
\label{section:experiments}
We conduct a series of controlled experiments specifically designed to evaluate {\textsc BLT} in relation to models utilizing tokenization, taking care to ensure that {\textsc BLT} does not benefit from access to longer sequence contexts that could yield an unfair advantage.

\subsection{Pre-training Datasets} In all our model scale experiments, we utilize two main pre-training datasets: (1) The Llama 2 dataset~\citep{touvron2023llama}, containing 2 trillion tokens sourced from a diverse range of publicly accessible materials, which undergo rigorous cleaning and filtering processes to enhance data quality; and (2) {\textsc BLT}-1T, a newly curated dataset consisting of 1 trillion tokens drawn from multiple public sources, which also integrates a segment of the pre-training data made available by Datacomp-LM~\citep{li2024datacomp}. The Llama 2 dataset serves as the foundation for scaling law studies that determine the ideal number of training tokens, following the methodology outlined by~\cite{dubey2024llama} for optimizing the {\textsc BLT} architecture. In contrast, the {\textsc BLT}-1T dataset is employed for an exhaustive pre-training regimen to enable direct comparison with Llama 3 on downstream tasks. Both datasets are explicitly devoid of any information harvested from Meta products or services. For tokenizer-based model baselines, we employ the Llama 3 tokenizer, which features a 128K token vocabulary. This tokenizer demonstrated superior baseline performance relative to the Llama 2 tokenizer in our assessments.

\subsection{Entropy Model} In our experiments, the entropy model is implemented as a byte-level language model, trained using the same data distribution as the overall {\textsc BLT} model. By default, this model is configured as a transformer comprising 100M parameters, 14 layers, a hidden size of 512, and a sliding window attention mechanism that spans 512 bytes. The rest of the hyperparameters match those used in both our local and global transformer setups. We also varied the size of the models, receptive field lengths, and underlying architectures to conduct further analysis, as detailed in \autoref{section:ablations}. Notably, for models with sufficiently restricted receptive fields, it becomes feasible to represent the entropy model in the form of a compact and efficient lookup table.

\subsection{Entropy Threshold and Equalizing Context Length} For entropy-based patching models, we determine a patching threshold that results in a target average \textit{patch size} computed over the pretraining data mixture. In {\textsc BLT}, as opposed to tokenization approaches, the \textit{patch size} is a flexible parameter and can be set freely, with this choice significantly affecting the model's effective context window. To ensure consistent average context length and to prevent models utilizing larger patch sizes from benefitting disproportionately, we enforce that the expected total number of bytes per batch is fixed. Consequently, for models employing larger patch sizes, we correspondingly decrease their sequence lengths. Specifically, for Llama 2 data, we allocate a context of 8k bytes, whereas for the {\textsc BLT}-1T dataset, we expand the average context window to 16k bytes, all while keeping the average batch size steady at 16M bytes.

Although the average batch size remains unchanged, dynamic patching strategies result in varying proportions of bytes per patch when constructing data batches. In our {\textsc BLT} training setup, we optimize performance by organizing batches to contain an equal number of patches, thereby eliminating the need for additional padding in the computationally intensive latent transformer. To prevent memory usage from surging due to exceptionally large patches, we pad and, if necessary, truncate input byte sequences to fixed lengths—12k for Llama 2 and 24k for the {\textsc BLT}-1T dataset—during training.

\subsection{Entropy Model Context}
\label{section:entropy-context}
Our empirical observations indicate that implementing entropy patching leads to increasingly larger patches, particularly in highly structured tasks such as multiple choice problems, which tend to display significant repetition (see the patching procedure on an MMLU sample in \autoref{fig:mmlu_patching}). These disparities stem from reduced entropy associated with the repeated elements in the entropy model context. Therefore, for the broad evaluation of {\textsc BLT}-Entropy at a patch size of 4.5, we refresh the entropy context using newline characters and adopt an approximate monotonicity constraint, as this approach is less susceptible to "entropy drift" caused by fluctuations in context length. Notably, this modification is limited to the method for computing entropies; our strategy for determining the appropriate entropy threshold remains unchanged.

\subsection{FLOPs Estimation} 
\label{section:flops}

Our methodology for calculating transformer flop counts is primarily based on the equations presented in Chinchilla~\citep{hoffmann2022training}, which detail the flop contributions from feed-forward networks, qkvo projections within self-attention, as well as the computation associated with attention and the output projection. However, our approach diverges by treating the input embedding layer as an optimized lookup rather than performing dense matrix multiplication, effectively attributing zero flops to this operation. In line with earlier studies, we approximate that the backward pass entails twice the flop count of the forward computation.

To determine the number of flops \textit{per byte} in {\textsc BLT} architectures, we aggregate the floating-point operations needed for the local encoder transformer, the global latent transformer, and the local decoder transformer, incorporating the computational costs from cross attention modules in both the encoder and the decoder:

\begin{align}
    \text{FL}_{\text{{\textsc BLT}}} &= \text{Transf. FL}(h_{\mathcal{G}}, l_{\mathcal{G}}, m=n_{ctx}/n_p,V=0) / n_p\\
   &+\text{Transf. FL}(h_{\mathcal{E}}, l_{\mathcal{E}}, m=w_{\mathcal{E}}, V=0) \\
   &+ \text{Transf. FL}(h_{\mathcal{D}}, l_{\mathcal{D}}, m=w_{\mathcal{D}}, V=256) \\ 
    &+ \text{Cross Attn. FL}(h_{\mathcal{E}}, l_{\mathcal{E}}, m=n_p, r=n_p/k) \times k/n_p \\ 
   &+ \text{Cross Attn. FL}(h_{\mathcal{D}}, l_{\mathcal{D}}, m=k, r=k/n_p) 
\end{align}

Here, $n_{ctx}$ denotes the sequence length measured in bytes, and $n_p$ indicates the patch size. The parameter $r$ signifies the proportion of queries to keys/values, while $k$ represents the ratio of the patch dimension to the byte dimension, effectively describing how many localized model segments are combined to generate the overall model embedding ($k=2$ as shown in \autoref{fig:crossattn}). The variable $V$ signifies the size of the vocabulary used for the output layer, relevant exclusively for the local decoder. The type of sequence (byte-level or patch-level) on which a module is utilized determines the context length $m$ employed in the attention mechanism. We adapt the computation of attention flops in alignment with the characteristics of each module. Precise formulas detailing how to compute the flops for both Transformer and Cross-Attention mechanisms are available in Appendix~\ref{section:flops-appendix}.

\subsection{Bits-Per-Byte Estimation} Perplexity serves as a measure of prediction uncertainty per token and is only meaningful under a consistent tokenizer framework. For evaluations spanning both token-level and byte-level models, we adopt the Bits-Per-Byte (BPB) metric—consistent with earlier research~\citep{xue2022byt5, yu2023megabyte, wang2024mambabyte}—which provides a tokenizer-agnostic analogue to perplexity. More precisely:

\begin{align}
    \text{BPB}(x)&= \frac{\mathcal{L}_{CE}(\pmb{x})}{\ln(2)\cdot n_{\text{bytes}}}
\end{align}

where the uncertainty over the data $\pmb{x}$ as measured by the sum of the cross-entropy loss is normalized by the total number of bytes in $\pmb{x}$ and a constant. 

\subsection{Transformer Architecture Hyperparameters} 
For each transformer block within {\textsc BLT}, encompassing both local and global model components, our configuration is principally aligned with Llama~3~\citep{dubey2024llama}. Specifically, the feed-forward networks incorporate the SwiGLU activation function~\citep{swiglu}, while self-attention layers are augmented with rotary positional embeddings (RoPE)~\citep{RoPe} adopting $\theta=500000$~\citep{xiong2024effective}. For normalization within layers, we employ RMSNorm~\citep{RMSNorm}. Flash attention~\citep{dao2022flashattention} is utilized across all self-attention layers that operate using standard attention masks, including both \textit{block causal} and \textit{fixed-window block causal} configurations, with a window size set to 512 for the latter. Given that cross-attention modules necessitate dynamically generated patch-specific masks, we leverage Flex Attention\footnote{\url{https://pytorch.org/blog/flexattention}}, which allows for fused kernel execution and thus markedly accelerates training.

\subsection{{\textsc BLT}-Specific Hyperparameters}

To evaluate the performance of the {\textsc BLT} models, our experimental framework encompasses two major aspects: analysis of scaling behavior and assessments on downstream tasks. The set of model sizes explored includes 400M, 1B, 2B, 4B, and 8B parameters; detailed architectural hyperparameters for each can be found in Appendix~Table~\ref{tab:arch}.

For initializing queries in the first cross-attention stage of the local encoder, we utilize max-pooling. All {\textsc BLT} variants use $500,000$ hashes employing a single hash function, with n-grams varying from size 3 to 8. Consistent across models, the learning rate is set at $4e-4$. This fixed learning rate for both the token-based and {\textsc BLT} models is justified by preliminary hyperparameter tuning within the $[1e-3, 1e-4]$ range, where results at the 400M and 1B scales indicate $4e-4$ as optimal. In experiments tracking scaling on Llama-2 datasets, training batch-sizes are set according to recommendations by~\cite{dubey2024llama} or a byte-wise equivalent.

For training, we employ the AdamW optimizer~\citep{adamw}, configuring $\beta_1$ at 0.9, $\beta_2$ at 0.95, and setting $\epsilon=10^{-8}$. The learning rate undergoes 2000 steps of linear warm-up before transitioning to a cosine annealing schedule down to zero. A weight decay of 0.1 is applied, along with global gradient norm clipping at 1.0.

\section{Scaling Trends}
\label{section:scaling}

We provide a comprehensive overview of the scaling behavior observed in byte-level models, offering insights valuable for guiding future scaling of {\textsc BLT} models. Our analysis extends previous investigations of byte-level modeling by: (a) examining performance trends under compute-optimal training configurations; (b) training comparable 8B parameter models with substantial data volumes (as much as 1T tokens or 4T bytes) and evaluating these models on various downstream applications; and (c) analyzing scaling patterns under scenarios where inference cost is kept constant. Subsequent sections will delve into particular benefits associated with modeling sequences at the byte level.

\subsection{Parameter Matched Compute Optimal Scaling Trends}
\label{section:parameter-matched-scaling}
With the Llama 2 dataset, we conduct training of a series of \textit{compute-optimal} bpe and {\textsc BLT} models at four distinct parameter scales, from 1B to 8B parameters. We subsequently examine the relationship between training flops and language modeling accuracy over a representative portion of the training data mixture. For bpe models, the allocation of parameters to dataset size adheres to the compute-optimal ratio outlined in Llama~3~\citep{dubey2024llama}. This approach to compute-optimality is intended, from a theoretical standpoint, to produce maximal training set performance under a specified computational constraint~\citep{hoffmann2022training}, thereby establishing a solid baseline for comparison. Parallel to each bpe model, a corresponding {\textsc BLT} variant is trained using the same dataset, and employing a Latent Transformer with an identical configuration and parameter count as its bpe Transformer counterpart.

As depicted in~\autoref{fig:scaling-compute-opt} (right), {\textsc BLT} models consistently equal or surpass the performance of bpe models, and this pattern persists as model size and FLOPs increase. To our knowledge, {\textsc BLT} represents the first byte-level Transformer architecture to exhibit comparable scaling behavior to BPE-based models within compute-optimal settings. These findings thus support our hypothesis that the optimal parameter-to-training-compute ratio observed for bpe models is similarly applicable to {\textsc BLT}, or at least does not deviate significantly.

Both enhancements to architecture and the introduction of dynamic patching are essential for aligning with BPE scaling behavior. In ~\autoref{fig:scaling-compute-opt} (left), we present a comparison between models employing space-patching and Llama 3. To estimate the performance of SpaceByte \citep{slagle2024spacebyte}, we use the {\textsc BLT} approach to space-patching, omitting both n-gram embeddings and cross-attention mechanisms. While SpaceByte offers gains relative to Megabyte, a considerable performance gap still separates it from Llama 3. The right panel of \autoref{fig:scaling-compute-opt} highlights the impact of combining architectural refinement with dynamic patching; at large scale, {\textsc BLT} models achieve competitive results with leading tokenizer-based models such as Llama 3.

In addition, our findings indicate that the specific tokenizer utilized influences performance in tokenizer-based models: employing the Llama-3 tokenizer yields better results than the Llama-2 tokenizer, even when models are trained on identical datasets.

In summary, when employing much larger patch sizes, our {\textsc BLT} framework demonstrates intermediate scaling behavior relative to Llama 2 and Llama 3. The average token lengths for Llama 2 and Llama 3, based on their bpe tokenizers, are 3.7 and 4.4 bytes, respectively. By comparison, {\textsc BLT} maintains analogous scaling dynamics while utilizing mean patch sizes of 6 or even 8 bytes. Since inference flops decrease with larger average patch sizes, a configuration with 8-byte patches could result in close to a 50\% reduction in inference flops. Furthermore, scaling both model and dataset sizes tends to benefit models that use increased patch sizes. Notably, although {\textsc BLT} with 8-byte patches initially underperforms relative to bpe-tokenized Llama 2 at the 1B scale, it surpasses the bpe baseline by the 7B scale. This pattern points toward the possibility that adopting even larger patch sizes could be advantageous as model capacity and available training resources continue to increase.

\subsection{Beyond Compute Optimal Task Evaluations}
\label{section:evals}
To further investigate scaling characteristics, we extend training of an 8B {\textsc BLT} model past the compute-optimal limit, utilizing the {\textsc BLT}-1T dataset, which is both larger and of higher quality. Model performance is then assessed on a diverse set of established classification and generative benchmarks. For these evaluations, we target benchmarks encompassing common sense reasoning, factual knowledge, and code generation capabilities:
\paragraph{Classification tasks} cover ARC-Easy (0-shot)~\citep{Eval_ARC}, Arc-Challenge (0-shot)~\citep{Eval_ARC}, HellaSwag (0-shot)~\citep{Eval_hellaswag}, PIQA (0-shot)~\citep{Eval_piqa}, and MMLU (5-shot)~\citep{hendrycks2020measuring}. We apply a prompt scoring approach, which involves computing the probability across the possible choice tokens and summarize outcomes using mean accuracy scores.
\paragraph{Coding related generation tasks:} For code generation capabilities, we present pass@1 metrics for MBPP (3-shot)~\citep{Eval_MBPP} and HumanEval (0-shot)~\citep{Eval_HumanEval}, measuring the proficiency of large language models at Python code synthesis.

In \autoref{tab:evals}, we present a comparison among three models trained using the {\textsc BLT}-1T dataset: a model based on the Llama 3 tokenizer with BPE,\footnote{We opt for the Llama 3 tokenizer and its 128k vocabulary, as it delivers improved results over the Llama 2 variant, which utilizes a 32k vocabulary.} along with two {\textsc BLT} model variants. The first variant adopts a space-patching method ({\textsc BLT}-Space), while the second employs a patching strategy driven by entropy measurements ({\textsc BLT}-Entropy). The latter also incorporates an approximate monotonicity constraint and resets its context at new lines, consistent with the procedure described in~\autoref{section:entropy-context}. Training for all three models is conducted under matched computational (FLOP) resources. Notably, for {\textsc BLT}-Entropy, we further adjust the entropy threshold at inference, lowering it from 0.6 to 0.1; this modification enhances performance on evaluated tasks but increases the number of inference steps required.

The {\textsc BLT}-Entropy model achieves superior results compared to the Llama 3 model on four of the seven evaluated tasks, even though both models are trained with an equivalent amount of data. These gains can likely be attributed to two main factors: (1) more effective utilization of computational resources during training through dynamic patching, and (2) the model’s focus on directly representing byte-level data rather than token-level representations.

In contrast, {\textsc BLT}-Space lags behind the Llama 3 tokenizer in performance across nearly all tasks, except for one. However, its use of a larger mean patch size of 6 bytes leads to a marked decrease in inference flops. For comparison, both bpe and models utilizing entropy-patching maintain an average patch size of around 4.5 bytes on the combined training dataset. When constrained to an identical training compute budget, the model with the larger patch size processes 30\% more data than the other two models, which could move {\textsc BLT} even further from achieving compute-optimality.

\subsection{Patches Scale Better Than Tokens} The {\textsc BLT} framework enables concurrent scaling of both the model and patch sizes without exceeding the original training and inference computational costs or requiring more training data. This flexibility to enlarge patch size distinguishes patch-based architectures from traditional fixed-vocabulary token-based models, allowing them to overcome efficiency constraints, as described in Section~\ref{section:bpe-patch}. Increasing patch length reduces computational demand, freeing up resources to expand the global latent transformer, which processes information less frequently as a result.

We perform a fixed inference scaling experiment to examine whether increasing model size while reducing the number of processing steps on larger patches yields superior performance compared to smaller models executing more steps. Using base models of 400m and 3.6B parameters tokenized with Llama 2, we identify flop-equivalent counterparts utilizing the Llama 3 tokenizer as well as {\textsc BLT}-Entropy models with mean patch sizes of 6 and 8 bytes across the training datamix (refer to~\autoref{tab:fixed-inf-params} for detailed model specifications). For models employing a patch size of 8, we opt for 3 encoder layers rather than a single layer. All models are trained under different total flop constraints.

\autoref{fig:fixed_inference_scaling} demonstrates that {\textsc BLT} models exhibit more favorable scaling behavior compared to tokenization-based models across both inference flop categories. For both, BPE models display superior performance when training resources are limited; however, as training compute increases—shortly past the compute-optimal threshold—{\textsc BLT} models overtake them. In real-world scenarios, allocating additional resources to pretraining, even as a one-time cost, may be advantageous if it results in a higher-performing model under a constant inference budget. This is exemplified by 8B parameter models such as Llama 3.1, which have been pretrained on data volumes two orders of magnitude greater than what would be theoretically optimal given their size.

The point at which {\textsc BLT} surpasses token-based approaches has moved nearer to the compute-optimal regime as we advance to models with greater flop capacities (shifting from approximately 3x to 2.5x above the compute-optimal allocation). Furthermore, the patch size 8 variant demonstrates a more pronounced scaling behavior in the larger flop category, overtaking alternative models at an earlier stage. Referencing Section~\ref{section:parameter-matched-scaling}, it can be observed that larger patch sizes yield performance increasingly aligned with BPE models as model sizes grow. We partly explain this phenomenon by noting the proportionally reduced computational share consumed by the byte-level Encoder and Decoder components, whose scaling rate is outpaced by that of the Latent Transformer. For instance, an expansion of total parameters by 20-fold—from 400M to 8B—results in only about a twofold increase in {\textsc BLT}'s localized model parameters. This relationship holds significance because the introduction of larger patch sizes influences only the patch Latent Transformer’s computational cost, sparing the byte-level elements. Notably, this helps to clarify why, for the {\textsc BLT}-Entropy model with ps=8, the parameter ratio compared to the Llama 2 baseline marginally increased from 1.6x to 1.7x when scaled up to larger model sizes.

In conclusion, the results of our patch-length scaling analysis indicate that the {\textsc BLT} patch-based framework exhibits superior scaling behavior when both patch size and model size are jointly scaled up. These favorable scaling effects are maintained, and potentially enhanced, as model scale increases.

\section{Byte Modeling Improves Robustness} Additionally, we evaluate the robustness of {\textsc BLT} against token-based counterparts that do not utilize byte-level representations, and we introduce a method for augmenting existing token-based pretrained models with byte-level capabilities.
\label{sec:robustness}

\subsection{Character-Level Tasks}

An initial impetus for developing byte-level models stemmed from their inherent resilience to noise at the byte level within inputs, as well as their capacity to recognize the subcomponents of tokens—a task where models relying on traditional tokenizers often encounter difficulties. In order to assess these capabilities, we conduct supplementary experiments on benchmarks designed to test both the robustness against input noise and the sensitivity to character-level details, across English and multiple languages, encompassing not only alphabetic characters but also digits and phonemes. The outcomes of these experiments are summarized in Table~\ref{tab:char_tasks}.

\paragraph{Noisy Data} To assess the robustness of {\textsc BLT} relative to tokenizer-based architectures, we construct perturbed variants of the benchmark classification datasets introduced in Section~\ref{section:evals}. Specifically, we implement five unique character-level noise augmentation techniques: (a) \textit{AntSpeak}: The entire sequence is represented in uppercase characters, separated by spaces. (b) \textit{Drop}: Ten percent of the characters in the input are randomly omitted. (c) \textit{RandomCase}: Half of the characters in the text are converted to uppercase, while the remaining half are randomly set to lowercase. (d) \textit{Repeat}: Twenty percent of characters are redundantly repeated, up to four instances. (e) \textit{UpperCase}: All characters are uniformly converted to uppercase. For evaluation, each noising procedure is independently applied to the prompt, completion, or both, and the resulting scores are averaged. The outcomes, summarized in Table \ref{tab:char_tasks} for the noised HellaSwag~\citep{Eval_hellaswag} dataset, reveal that {\textsc BLT} consistently exceeds the robustness of tokenizer-based models, offering an average performance increase of 8 points over a model with identical training data. Additionally, {\textsc BLT} surpasses the Llama 3.1 model, despite the latter's training on a substantially larger corpus.

\paragraph{Phonology - Grapheme-to-Phoneme (G2P)} We evaluate how effectively {\textsc BLT} converts sequences of graphemes—that is, the characters in a word—into their corresponding phonemic transcriptions, which represent the word’s pronunciation. Table \ref{tab:char_tasks} displays the outcomes for the G2P task conducted in a 5-shot scenario utilizing Phonology Bench~\citep{suvarna2024phonologybench}. The results indicate that {\textsc BLT} achieves higher performance compared to the baseline Llama 3 1T tokenizer-based model.

\paragraph{CUTE}
To evaluate capabilities at the character level, we tested {\textsc BLT} using the CUTE benchmark~\citep{edman2024cute}, which encompasses multiple tasks grouped into three main areas: compositional understanding, recognition of orthographic similarity, and sequence manipulation. This benchmark is particularly demanding for most models reliant on tokenizers, as such models often have some knowledge about the spellings of their tokens but face difficulties effectively leveraging this information for text manipulation. Results in Table~\ref{tab:char_tasks} reveal that {\textsc BLT}-Entropy exceeds both BPE Llama 3 variants by over 25 points on the benchmark. Notably, our model exhibits outstanding capability on character manipulation tasks, achieving 99.9\% accuracy on both spelling-related tasks. These substantial gains, even though {\textsc BLT} was trained on 16 times less data compared to Llama 3.1, highlight the challenges BPE-based models encounter when acquiring character-level knowledge. Figure \ref{fig:cute_examples} presents examples demonstrating instances where the Llama 3 tokenizer-based model falters, whereas our {\textsc BLT} model succeeds. The only exceptions are the word deletion and insertion tasks, where BPE slightly outperforms; nonetheless, the margin is small, and the progression from character-level modeling to word-level manipulation may present fewer difficulties than the reverse. The evaluation protocols and original Huggingface prompts remain consistent across all tasks. Further prompt tuning could potentially enhance the performance of BPE models.

\paragraph{Low Resource Machine Translation} We assess the performance of {\textsc BLT} for both directions of translation—into and out of English—covering six widely recognized language families and twenty-one low-resource languages that use a variety of scripts, as defined in the FLORES-101 benchmark~\citep{goyal2022flores}. SentencePiece BLEU results are presented in Table~\ref{tab:flores}. Our findings indicate that {\textsc BLT} surpasses a model based on the Llama 3 tokenizer, securing an overall improvement of 2 BLEU points in translations into English and a 0.5-point increase in translations originating from English. For frequently studied language pairs, {\textsc BLT} delivers performance on par with, or marginally better than, Llama 3. Notably, in many cases involving underrepresented language pairs from low-resource families, {\textsc BLT} demonstrates superior performance compared to Llama 3, highlighting the strengths of byte-level modeling for handling rare or unseen byte patterns.

\subsection{Training {\textsc BLT} Using Llama 3 Initialization}
\label{sec:distillation}
We investigate a methodology that enables {\textsc BLT} models to benefit from pre-existing pre-trained tokenizer-based architectures to improve training efficiency and speed up convergence. Specifically, we initialize the global transformer parameters of {\textsc BLT} with those from a pre-trained Llama 3.1. After initialization, the global transformer's parameters are further updated with a learning rate set at one-tenth that of the local encoder and decoder, and optimization proceeds for the standard number of Llama 3 training steps. A detailed comparison with a baseline {\textsc BLT} is provided in Table~\ref{tab:distill}. The results clearly show that {\textsc BLT} initialized from Llama 3.1 demonstrates substantial performance gains over both the Llama 3 and baseline {\textsc BLT} models when controlled for equivalent FLOP budgets. Furthermore, in comparison to our {\textsc BLT}-Entropy variant reported in Table~\ref{tab:evals}—which was trained on a much larger corpus (1T tokens)—the {\textsc BLT} model initialized from Llama 3.1 nonetheless exhibits superior results on the MMLU benchmark. This indicates that this initialization strategy can lead to meaningful reductions in the training FLOPs required for strong performance.

This approach may be interpreted as converting models reliant on tokenizers into ones that operate without them, thereby transforming a pre-trained LLaMA 3.1 into a {\textsc BLT} model. For a thorough evaluation, we present in Table \ref{tab:distill} both the original LLaMA 3.1 model, which was trained on 15T tokens, and assess it alongside the {\textsc BLT} version derived from LLaMA 3. Our findings show a minor decrease in performance on the MMLU and HumanEval benchmarks, whereas declines on additional tasks are more pronounced. These outcomes indicate that further research is necessary to better utilize the pre-trained foundation and enhance results, especially by refining data composition and tuning hyperparameters.

\section{Ablations and Discussion}
\label{section:ablations}
This section presents ablation studies that validate the architectural decisions underlying {\textsc BLT}, as well as the selection of patching methodologies and hyper-parameters for the {\textsc BLT} model with 8B parameters trained on the {\textsc BLT}-1T dataset.

\paragraph{Entropy Model Hyper-parameters} To evaluate how the size of the entropy model and the length of its context window influence performance at scale, we train a range of byte-level entropy transformer models whose parameter counts span from 1 million to 100 million, and experiment with context window lengths ranging between 64 and 512. The relationship between bits per byte (bpb) and training computational expenditure is visualized through scaling law curves, constructed using the $400m$ and $1b$ {\textsc BLT} models on the Llama-2 dataset, as depicted in \autoref{fig:entropy_ablation}. Our analysis reveals that while increasing the entropy model size and the context window both improve scaling performance, the incremental benefit tapers off after reaching 50 million parameters.

\paragraph{Types of Patching}

We conduct an ablation study on the four patching strategies outlined in Section~\ref{section:patching}: (1) Strided Patching employing strides of 4 and 6, (2) patching at whitespace boundaries, (3) patching determined by BPE Tokenizer splits as per the Llama 3 tokenizer, and (4) entropy-driven patching utilizing a compact byte-level LLM.

Although dynamic patching shortens the actual sequence lengths, we ensure that the contextual window remains comparable across all patching approaches by regulating sequence length. Consequently, each model is exposed to an equivalent expected byte count per sequence during both training and evaluation, thereby mitigating any confounds related to disparities in modeled context size. As depicted in Figure~\ref{fig:scaling-compute-opt}, these controlled ablations demonstrate that every alternative patching scheme surpasses static patching, with space patching performing nearly as well as dynamic entropy based patching.

In \autoref{tab:evals_ablation}, we report benchmark results for {\textsc BLT} models using three different patching approaches—tokenizer-based methods, space patching, and entropy-based patching—all trained for an optimal duration on the Llama 2 dataset~\citep{dubey2024llama}. While space patching offers a more straightforward method, as it bypasses the need to compute entropy in real time during training, our findings demonstrate that the improvements observed with entropy-based patching in scaling analyses (see Section~\ref{section:scaling}) extend to downstream evaluation tasks as well.\footnote{The space patching outcomes shown are from initial experiments conducted without cross-attention, but similar patterns emerge when cross-attention is included.}

\paragraph{Cross-Attention} 
\autoref{tab:crossattn_ablations_1b} presents an ablation study on inserting cross-attention modules at different locations within the encoder and decoder of {\textsc BLT}. In the encoder, we explore three methods for initializing the queries used in cross-attention: (1) employing an identical learned embedding for all global states, (2) utilizing a hash embedding generated from the patch bytes, and (3) aggregating the encoder layer's hidden representations corresponding to the patch bytes via pooling.

Our results indicate that incorporating cross-attention within the \textit{decoder} yields the best performance. Introducing cross-attention in the encoder offers modest gains, which are observed exclusively when the queries are initialized via pooling. Furthermore, cross-attention proves particularly beneficial on the Common-Crawl dataset and becomes even more advantageous as patch size increases.

\paragraph{n-gram Hash Embeddings} 

We experimented with n-gram hash embedding vocabulary sizes of 0, 100K, 200K, and 400K, and the corresponding outcomes are shown in Table~\ref{tab:ngram_ablations_1b}. Our analysis demonstrates that incorporating hash embeddings yields improvements across all evaluated domains, with the most pronounced benefits observed for Wikipedia and Github datasets (resulting in a 0.04 bpb improvement, as opposed to just 0.01 bpb after 15k steps at 8B scale). Reducing hash vocabularies from 500K to 300K at the 8B scale caused only a minor performance decrease of approximately 0.001 bpb after 15k steps. These results suggest that hash embeddings are crucial for enabling {\textsc BLT} to achieve performance on par with models that utilize tokenizers; nonetheless, beyond 300K hashes, performance improvements become marginal. Furthermore, evidence suggests that the advantages conferred by hash embeddings and cross-attention are largely independent, as they offer enhancements on distinct datasets.

\paragraph{Local Model Hyperparameters} Table \ref{tab:local_layers} presents an ablation study examining different configurations for the number of layers in both the local encoder and decoder. Results indicate that, when hash n-gram embeddings are utilized, {\textsc BLT} achieves strong performance with a very shallow encoder—comprising only a single layer—while benefiting from a comparatively deeper decoder.

\section{Related Work}
\label{section:related_work}

\textbf{Character-Level RNNs:} The task of Character Language Modeling has been prominent since the advent of neural network approaches~\citep{sutskever2011generating,mikolov2012subword,graves2013generating}, due in large part to their inherent capacity to handle out-of-vocabulary terms natively without the need for fallback strategies. \cite{kim2016character} introduced a model that focuses exclusively on character sequences as input, utilizing convolutional and highway networks as input encoders to LSTM-based RNNs. This approach achieved comparable performance to then-current RNN-based language models on English and delivered superior results for morphologically complex languages, further highlighting the appeal of character-level LLMs in these contexts. \cite{kenter2018byte} applied byte-level LSTM architectures to machine comprehension, surpassing traditional word-level models especially in agglutinative languages such as Turkish and Russian. Extending this line of inquiry, \cite{zhang2015character} employed convolutional networks over characters for classification purposes, demonstrating improved outcomes over word-based techniques for certain scenarios. Hierarchical LSTM models utilizing boundary detectors to extract latent textual hierarchies were proposed by \citet{chung2019hierarchical}, providing additional performance gains in character-level language modeling. Diverging from attention mechanisms, ByteNet~\citep{kalchbrenner2016neural} employed convolutional layers at the character level to address machine translation tasks.

\textbf{Character-Level Transformers:} The advent of transformer architectures leveraging attention mechanisms~\citep{vaswani2017attention} and subword tokenization techniques~\citep{sennrich-etal-2016-neural} led to considerable enhancements in neural language models' effectiveness on a variety of tasks and standard benchmarks. Nonetheless, relying on word or sub-word units introduces an inherent inductive bias about the abstraction level at which these models function. Seeking to bridge the achievements of transformer models with encouraging outcomes observed in character-level language modeling, \citet{al2019character} constructed deep transformer networks and, employing auxiliary loss functions, trained models that surpassed earlier character-level LSTM-based language models. Despite these advancements, their models continued to lag considerably behind large language models employing word-level representations. Similarly, GPT-2~\citep{radfordgpt2} found that when dealing with expansive datasets such as the 1 billion word benchmark, language models using byte-level representations did not achieve performance on par with their word-level counterparts.

Although \cite{Choe2019BridgingTG} showed that transformer-based byte-level LLMs can surpass subword-level LLMs of similar scale, these models demand considerably higher computational resources and training time. Similarly, \cite{el2020characterbert} introduced CharFormer, a BERT variant that constructs word representations through convolutions on character embeddings, noting performance gains in the medical domain, albeit with substantially greater computational cost. In their work on CANINE, \cite{clark2022canine} presented an encoder-only model with 150M parameters designed to operate over raw character sequences. CANINE’s architecture integrates a deep transformer—conceptually parallel to our global model—as well as a combination of strided convolutions and a local transformer component for downsampling input characters. This configuration yields superior performance compared to token-level encoder-only architectures such as mBERT on a range of multilingual downstream tasks. The ByT5 model~\citep{xue2022byt5} examined encoder-decoder frameworks at the byte level that avoid patching operations altogether. Despite ByT5 showing resilience to input noise and performing on par with models using tokenization while being trained on a quarter of the data, the absence of patching led to computational bottlenecks, as the self-attention mechanism must operate across every byte, escalating computational demands significantly. Modelling data at the byte level, rather than with subword units, lengthens sequences and therefore presents substantial scalability challenges for these models. In more recent advances, employing the Mamba Architecture~\citep{gu2023mamba}, which supports a fixed-size memory state even across long contexts, \cite{wang2024mambabyte} have developed a byte-level Mamba-based model—again foregoing patching—which outperforms transformer-based byte-level models on several datasets (in terms of bits-per-byte), specifically at the 350M parameter scale within a computationally controlled setting.

\textbf{Patching-based approaches:} Leveraging patching techniques can significantly decrease the computational load, particularly the number of floating point operations (flops), required for byte-level LLMs, while still maintaining strong performance. Numerous studies have shown that, at modest model sizes and with limited training bytes, patching proves effective. For example, \cite{nawrot-etal-2022-hierarchical} utilize static patching, incorporating downsampling and upsampling within their hourglass transformer architecture, which achieves superior results compared to other byte-level baselines at the 150M parameter scale. Building on this, \cite{nawrot-etal-2023-efficient} introduce more advanced dynamic patching strategies, such as an end-to-end learned boundary predictor, supervision through specific tokenizers, and an entropy-driven patching approach akin to {\textsc BLT}. Their results indicate performance gains over standard transformers on language modeling benchmarks using 40M parameter models and approximately 400M tokens. Further, \cite{lester2024training} analyze training on text sequences compressed with arithmetic coding, which enables greater compression ratios than traditional BPE. Employing a technique called equal-info windows, they report improvements over byte-level baselines on language modeling tasks, albeit without surpassing models trained with subword units.

Our research is primarily motivated by, and closely aligned with, MegaByte~\citep{yu2023megabyte}, a causal decoder-only large language model (LLM) that employs a fixed, static patching scheme, concatenating byte representations to form patches, which are then translated back to bytes on the decoder end via a local model. Their findings indicate that MegaByte attains performance comparable to that of tokenizer-based approaches at the 1B parameter scale when evaluated on a dataset containing approximately 400B bytes. Across all our experimental settings, we include MegaByte as a baseline and observe that static patching still trails behind state-of-the-art compute-efficient tokenizer-based models when flops are held constant; our work further illustrates how {\textsc BLT} helps to close this performance gap. Similar observations are made by \cite{slagle2024spacebyte}, who propose expanding the static patching strategy to include whitespaces and similar byte classes, in addition to introducing a local encoder model. Their experiments reveal enhancements over token-based transformer architectures under compute-constrained settings in domains such as Github and arXiv for models with 1B parameters. We conduct evaluations using their model as well, concluding that even more significant architectural innovations will be essential for byte-level models to scale effectively and genuinely match leading tokenizer-based systems like Llama 3.

\section{Limitations and Future Work}

For the selection of architectural configurations in this study, we train our models using the optimal number of steps identified for Llama~3~\citep{dubey2024llama}. Nevertheless, as these scaling laws were originally derived for BPE-level transformer models, their direct application may result in less-than-ideal ratios between data and parameters for the {\textsc BLT} framework. The derivation of suitable scaling laws specific to {\textsc BLT} is reserved for future investigations, which could reveal even more advantageous scaling behaviors in our model design. Furthermore, since many of our experimental trials were limited to model sizes up to 1B parameters, there remains the possibility that the most effective architectural settings may shift when scaling up to 8B parameters or higher, potentially enabling superior performance at larger model sizes.

Current transformer frameworks and repositories are optimized primarily for token-based transformer models. Although we conduct experiments with theoretical FLOP equivalence and leverage specialized implementations—like FlexAttention—for managing architectural modifications that diverge from standard transformers, our versions might not match tokenizer-based systems in actual runtime performance and could still require additional optimization efforts to reach comparable efficiency.

Whereas {\textsc BLT} employs an independently trained entropy model for the patching process, integrating the learning of the patching model into an end-to-end training scheme represents a promising avenue for future research. As described in Section \ref{sec:distillation}, we conduct preliminary studies that demonstrate early evidence of effective "byte-ification" of tokenizer-based architectures like Llama 3, which have been trained on datasets exceeding 10T tokens, by reusing and freezing their global transformer parameters. Pursuing this line of research further may yield techniques that both preserve the advantages introduced by bytefication and even enhance performance relative to these tokenizer-based models, all without necessitating training from scratch.

\section{Conclusion}
\label{section:conclusion}

In this work, we introduce the Byte Latent Transformer (\textbf{BLT}), an architectural innovation that challenges the standard reliance on fixed-vocabulary tokenization commonly used in large language models. Through a flexible, learnable strategy for clustering bytes into patches, {\textsc BLT} dynamically distributes computation in proportion to data intricacy, thereby boosting both robustness and computational efficiency. Our comprehensive scaling analysis reveals that {\textsc BLT} attains parity with established tokenization-driven models such as Llama 3 across scales up to 8B parameters and 4T bytes, while also allowing for reductions of up to 50\% in inference flops for a marginal drop in evaluation scores. In addition, {\textsc BLT} provides an avenue to simultaneously expand both the model size and patch size without increasing inference cost, introducing new possibilities for scaling in real-world compute environments. Operating directly on byte-level inputs, {\textsc BLT} advances the handling of infrequent data distributions, enhances resilience to noisy inputs, and enables richer modeling of sub-word elements. Collectively, these findings establish {\textsc BLT} as a compelling alternative to conventional tokenization-based systems, presenting an efficient, scalable, and robust framework for the next generation of adaptable language models.

\end{document}